\section{Fase 2: Evolução Planejada da Modelagem (M8)}
\label{sec:m8-fase2}

As limitações declaradas na Seção~\ref{sec:limitacoes} não são acidentes, mas fronteiras
assumidas. A segunda fase do projeto (milestone~M8) ataca três delas com decisões registradas
nos ADR-0009, 0010 e 0011. A lógica reusável das três já está implementada em \texttt{src/}; a
agregação do erro (ADR-0009) foi, além disso, \textbf{calibrada e medida}
(Seção~\ref{sec:m8-pooling-result}). As demais aguardam os notebooks de orquestração. Esta
seção documenta o plano, a evidência disponível e a ordem de prioridade.

\subsection{Agregação do erro por janela: \emph{max}/percentil (ADR-0009)}
\label{sec:m8-pooling}
Hoje o erro de uma janela é a \textbf{média} do erro absoluto sobre os 30 passos
(Seção~\ref{sec:m4}). Um choque sintético de um único dia entra com peso $1/30$ e é
\emph{diluído} pelos 29 dias normais, podendo permanecer abaixo do limiar --- a causa direta do
teto de \emph{recall} em choques pontuais (Seção~\ref{sec:limitacoes}). A proposta substitui a
agregação por uma parametrizável: \emph{mean} (atual), \emph{max} (pior passo da janela) ou
\emph{percentil} alto (p.\,ex.\ p90, meio-termo robusto a um único ponto ruidoso). A redução é
feita em duas etapas e nesta ordem: média sobre o eixo de \emph{features}, depois agregação
sobre o eixo \emph{temporal}.
\par\noindent\textbf{Cuidado metodológico:} \emph{max}/percentil alteram toda a distribuição do
erro; o limiar p95, calibrado sobre a média, \textbf{não} transfere e precisa ser recalculado
sobre o mesmo escore agregado do treino (sob pena de inflar falsos positivos).

\subsubsection*{Evidência empírica (notebook \texttt{07\_aggregation\_recalibration})}
\label{sec:m8-pooling-result}
O limiar estático foi recalibrado por agregação e a injeção sintética de M5 ($k\sigma$) foi
reavaliada nos quatro ativos. O p95 de \emph{max}/percentil é $\approx$3--4$\times$ o da média
(p.\,ex.\ PETR4: $0{,}124$ vs.\ $0{,}408$), confirmando que reaproveitar o limiar antigo seria
um erro. A Tabela~\ref{tab:m8-pooling} resume as métricas médias (limiar estático recalibrado):

\begin{table}[h]
\centering
\caption{Agregação do erro por janela vs.\ Precision/Recall/F1, média dos quatro ativos com
limiar estático recalibrado (ADR-0009).}
\label{tab:m8-pooling}
\begin{tabular}{lccc}
\toprule
\textbf{Agregação} & \textbf{Precision} & \textbf{Recall} & \textbf{F1} \\
\midrule
\texttt{mean} (anterior) & 0,548 & 0,162 & 0,229 \\
\textbf{\texttt{max}} & \textbf{0,842} & \textbf{0,351} & \textbf{0,468} \\
\texttt{percentile} (p90) & 0,706 & 0,209 & 0,288 \\
\bottomrule
\end{tabular}
\end{table}

\noindent O resultado é mais forte que o previsto: \emph{max} não apenas \textbf{dobrou o
\emph{recall}} ($0{,}162 \to 0{,}351$) como, contrariando a hipótese de um \emph{trade-off},
\textbf{também elevou a \emph{precision}} ($0{,}548 \to 0{,}842$). Com o choque de um dia
separando-se nitidamente da normalidade no pior passo e o limiar recalibrado, caem
simultaneamente falsos negativos e falsos positivos. Recomenda-se adotar \emph{max} na detecção;
o \emph{default} de configuração permanece \texttt{mean} até que M4--M6 sejam reportados sob a
nova agregação, preservando a reprodutibilidade dos números atuais.

\subsection{Validação Walk-Forward (ADR-0010)}
\label{sec:m8-walkforward}
A seleção de hiperparâmetros (notadamente \texttt{latent\_dim}) apoia-se hoje num \emph{holdout}
único, cuja estimativa tem alta variância. A proposta adota validação \emph{walk-forward} com
\texttt{TimeSeriesSplit}: fatias expansíveis do passado, avaliação no bloco futuro adjacente,
repetidas por $K$ folds, produzindo média~$\pm$~desvio do \texttt{val\_loss} por candidato. A
restrição metodológica central do projeto é preservada \textbf{por fold}: o \texttt{MinMaxScaler}
é reajustado apenas no treino de cada fold e as janelas são geradas \emph{dentro} de cada
partição, sem cruzar a fronteira treino/validação (ADR-0001). Esta mudança não altera o detector
em produção --- é instrumento de rigor e credibilidade na escolha do modelo final.

\subsection{Entrada multivariada OHLCV (ADR-0011)}
\label{sec:m8-multivariado}
O modelo é univariado (log-retorno de \texttt{Close}), tornando anomalias puramente de volume
invisíveis por construção (Seção~\ref{sec:limitacoes}). A proposta expande a entrada de
$(30,1)$ para um tensor multivariado, partindo do mínimo informativo --- \textbf{Close+Volume}
$(30,2)$ --- e só crescendo para OHLCV $(30,5)$ se houver ganho, dado que Open/High/Low/Close são
quase colineares. Riscos a controlar: escalas heterogêneas (volume $\sim 10^7$ vs.\ retorno
$\sim 10^{-2}$) exigem \texttt{MinMaxScaler} \textbf{por coluna}; a não-estacionariedade do volume
exige \texttt{log1p} antes de escalar (sob pena de saturação no range pós-2020); e o MSE
desbalanceado entre canais recomenda registrar o erro \textbf{por canal}, indispensável para
atribuir a anomalia (\emph{volume-driven} vs.\ \emph{price-driven}).

\subsection{Prioridade: maior retorno sobre o esforço}
\label{sec:m8-prioridade}
A ordem de execução é deliberada:
\begin{enumerate}[noitemsep]
  \item \textbf{Agregação \emph{max}/percentil (ADR-0009)} --- maior retorno por esforço:
        poucas linhas, ataca diretamente o \emph{recall} declarado, efeito imediato e mensurável.
  \item \textbf{Walk-Forward (ADR-0010)} --- não melhora o modelo, mas constrói a régua honesta
        para comparar a mudança anterior contra a \emph{baseline} e calibrar a próxima.
  \item \textbf{Multivariado OHLCV (ADR-0011)} --- maior potencial, maior superfície de risco;
        só compensa com (1) agregação que preserva o sinal e (2) validação confiável já no lugar,
        e ainda assim começando por Close+Volume.
\end{enumerate}
Em síntese: (1) é o ganho rápido de \emph{recall}, (2) é o microscópio para confiar nos números,
(3) é a aposta de maior alcance, viável apenas sobre as duas primeiras.
